\documentclass[10pt]{article}
\usepackage[a4paper,margin=25mm]{geometry}
\usepackage{amsmath,amssymb}
\usepackage{algorithm}
\usepackage{algpseudocode}
\usepackage{booktabs}
\usepackage[colorlinks=true,linkcolor=blue]{hyperref}

\newtheorem{remark}{Remark}

\title{\textbf{COSH: Closed-Form Onset Search via Hyperbolic Response}\\[2mm]
\large Methodology for pre-flight critical-moment identification of a
multirotor from ground-contact tip-over excitation}
\author{}
\date{}

\begin{document}
\maketitle
\vspace{-10mm}

\noindent
This document specifies the complete identification methodology, from raw
signals to the take-off feed-forward moment. Every stage mirrors the
implementation in \texttt{critical\_value\_getter\_piecewise.py}; the
companion analyses (model-fidelity bounds, excitation design, ground-effect
bracket) are documented in \texttt{fidelity\_section.tex} and are only
referenced here.

\section{Problem setting and notation}

A multirotor rests on its landing gear. Ramping the commanded roll (pitch)
moment tilts the vehicle until it begins to rotate about the landing-gear
contact line — the \emph{pivot}. The applied moment at that \emph{rotation
onset} is the critical moment $M_{\mathrm{crit}}$; comparing the two tip
directions yields the CoM offset and the feed-forward moment
$M_{\mathrm{ff}}$ that compensates take-off. No rig, no motion capture and
no stable in-flight data are used.

\begin{table}[h]\centering\small
\begin{tabular}{llll}
\toprule
Symbol & Meaning & Value & Provenance \\
\midrule
$W$, $z_{\mathrm{CoM}}$ & weight, CoM height & $31.59$~N, $\sim$0.3~m & measured / conservative \\
$l_{p,\phi},\,l_{p,\theta}$ & pivot offsets (roll, pitch) & $0.140$, $0.110$~m & conservative geometry \\
$J_P$ & inertia about the pivot line & — & never used numerically \\
$C_T$ & thrust map, $T_i = C_T\Omega_i^2$ & $1.3175\times10^{-7}$ & bench calibration \\
$M_{\mathrm{thr}}$ & window threshold & $0.01$~N\,m & noise floor \\
$M_{\mathrm{cap}}$ & allocator caps $(x,y)$ & $2.37$, $2.74$~N\,m & allocation feasibility \\
$\varepsilon_{\max}$ & ramp-slope gate & $3\%$ & sensitivity chain (Sec.~\ref{sec:gates}) \\
$N_{\min}$ & window-sample gate & $38$ & $30$ pre $+\,8$ post \\
$e_{\mathrm{lin,max}}$ & linearity-RMSE gate & $30$~mN\,m & $10\times$ noise floor \\
$N_{\mathrm{seg}}$ & sweep margin (each end) & $8$ & degenerate-candidate guard \\
\bottomrule
\end{tabular}
\caption{Design constants. $(C_2, K)$ are \emph{calibrated}, not set
(Sec.~\ref{sec:calibration}).}
\end{table}

\section{Contact-phase dynamics and the closed form}
\label{sec:closedform}

About the active pivot line, with the small-angle attitude dependence
retained ($\cos\varphi \approx 1$, $\sin\varphi \approx \varphi$), the
roll/pitch dynamics read
\begin{equation}
J_P\,\dot\omega
 = M(t) + f\,l_p - W\,(l_p + \lambda_{\mathrm{off}}) + W z_{\mathrm{CoM}}\,\varphi ,
\label{eq:dyn}
\end{equation}
where $M$ is the commanded body moment, $f$ the collective thrust and
$\lambda_{\mathrm{off}}$ the CoM offset along the excited axis. Below the
critical value the vehicle is static. At the onset the net moment
vanishes, which supplies \emph{two} anchoring conditions,
\begin{equation}
\omega(t_{\mathrm{crit}}) = 0, \qquad \dot\omega(t_{\mathrm{crit}}) = 0 .
\label{eq:onset-cond}
\end{equation}
Past the onset gravity feeds back positively: the linearized dynamics have
\emph{real} eigenvalues $\pm C_2$ and, with the measured moment ramp
$M(t) = M_{\mathrm{crit}} + \dot M\,\tau$ in onset-relative time
$\tau = t - t_{\mathrm{crit}}$, the exact solution of
\eqref{eq:dyn}--\eqref{eq:onset-cond} collapses to the hyperbolic
closed form
\begin{equation}
\boxed{\;\omega(\tau) = C_1\bigl(\cosh(C_2\tau) - 1\bigr), \qquad
C_1 = K\dot M,\quad K = \frac{1}{W z_{\mathrm{CoM}}},\quad
C_2 = \sqrt{\frac{W z_{\mathrm{CoM}}}{J_P}} \; }
\label{eq:cosh}
\end{equation}
for $\tau \ge 0$ (and $\omega = 0$ for $\tau < 0$). Because
$J_P C_2^2 = W z_{\mathrm{CoM}}$, the amplitude $C_1$ is \emph{fixed} by
the measured ramp rate: with $(C_2, K)$ pinned as rig constants and the
baseline fixed by continuity, \textbf{no shape parameter is free per run}.
The time-quadratic model of the original submission is the $C_2 \to 0$
limit of \eqref{eq:cosh} — it discards the gravity term that creates the
instability, with $14$--$75\%$ truncation error over the windows actually
used — and is retained only as a baseline.

\section{Signals and excitation window}
\label{sec:window}

Per run, the moment is reconstructed from the measured rotor speeds,
$M(t) = C_T \sum_i b_{a,i}\,\Omega_i^2(t)$, and $\omega(t)$ is taken from
the odometry rate. The window end is
$i_{\mathrm{end}} = \arg\max_t |M|$, \emph{truncated at the first sample
exceeding the allocator cap} $M_{\mathrm{cap}}$: past it the thrust
allocation saturates, the executed ramp is no longer linear, and the cosh
family does not apply. (On the reference dataset the cap never binds;
peak $|M| \le 1.7$~N\,m.) The window start is the \emph{last}
sub-threshold sample before the peak, walking backwards — robust to
spin-up blips that cross the threshold long before the actual ramp.

\section{Onset-free run gates}
\label{sec:gates}

Before any onset is estimated, each run must pass three gates evaluated on
the measured moment trace alone (hence free of circularity):
\begin{enumerate}
\item \textbf{Ramp-slope error} $|\varepsilon| \le 3\%$, with
$\varepsilon = (|\dot M| - \dot M_{\mathrm{cmd}})/\dot M_{\mathrm{cmd}}$
and $\dot M$ the least-squares slope over the window. A $3\%$ error in
$\dot M$ perturbs $C_1 = K\dot M$ by $3\%$ — about $7\times$ below the
$\pm20\%$ level the sensitivity analysis shows harmless.
\item \textbf{Window samples} $N_{\mathrm{full}} \ge 38$: a necessary
condition for any onset position. Because the window opens at
$|M| = 10$~mN\,m while $M_{\mathrm{crit}} \gtrsim 0.4$~N\,m, the onset
cannot occur earlier than
$\hat N_{\mathrm{pre}} = M_{\mathrm{crit}}^{\min}/(\dot M\,T_s) \ge 30$
samples into the window even at the fastest ramp; adding the
$N_{\mathrm{seg}} = 8$ post-onset fit minimum gives $38$.
\item \textbf{Linearity RMSE} $\le 30$~mN\,m about the run's own best-fit
line: a fault detector for aborted/stepped ramps, set at $10\times$ the
execution noise floor and above the healthy-run maximum.
\end{enumerate}
The gates run \emph{before} calibration so a poorly executed ramp cannot
contaminate the dataset-coupled constraint $C_1 = K\dot M$.

\section{Onset detection}
\label{sec:algorithm}

With $(C_2, K)$ pinned, detection reduces to an exhaustive sweep of the
single remaining unknown — the onset index. Each candidate is a plain
arithmetic evaluation: no optimizer, no initial guess, no seed model.

\begin{algorithm}[h]
\caption{COSH onset detection (single run)}
\label{alg:cosh}
\begin{algorithmic}[1]
\Require samples $\{(t_i,\omega_i)\}$ on the excitation window
  $[i_{\mathrm{start}}, i_{\mathrm{end}}]$; moment $M(t)$;
  calibrated $(C_2, K)$; measured ramp slope $\dot M$;
  margin $N_{\mathrm{seg}} = 8$
\Ensure onset index $j^\ast$; critical moment $M_{\mathrm{crit}}$
\State $C_1 \gets K\dot M$
  \Comment{onset conditions \eqref{eq:onset-cond} fix the amplitude}
\For{$j = i_{\mathrm{start}} + N_{\mathrm{seg}}, \dots,
  i_{\mathrm{end}} - N_{\mathrm{seg}}$}
  \Comment{exhaustive sweep of the single unknown}
  \State $C \gets \operatorname{median}\{\omega_i :
    i_{\mathrm{start}} \le i < j\}$
    \Comment{baseline by continuity}
  \State $\mathrm{cost}(j) \gets
    \sum_{i=i_{\mathrm{start}}}^{j-1}(\omega_i - C)^2
    + \sum_{i=j}^{i_{\mathrm{end}}}\bigl(\omega_i
    - [\,C_1(\cosh C_2 (t_i - t_j) - 1) + C\,]\bigr)^2$
\EndFor
\State $j^\ast \gets \arg\min_j \mathrm{cost}(j)$;\quad
  $M_{\mathrm{crit}} \gets M(t_{j^\ast})$
\end{algorithmic}
\end{algorithm}

Two implementation details. \emph{Baseline by median.} The pre-onset
gyro trace is heavy-tailed (propeller vibration, spin-up spikes), and —
more importantly — when a candidate $j$ overshoots the true onset, the
pre-segment tail contains the beginning of the hyperbolic rise. A mean
baseline drifts with that leaked rise and partially cancels the residual
it should expose, flattening the cost landscape; the median stays at the
rest level until half the segment is contaminated, so misplaced
candidates are penalized sharply and the minimum stays well defined. It
is also a one-shot, seedless statistic, consistent with the
no-optimizer design. An ablation rerunning the full pipeline — including
the $(C_2, K)$ calibration — with the mean instead
(\texttt{analysis/baseline\_stat\_ablation.py}) leaves the run gating
unchanged and $62/132$ onsets identical, but biases the remainder toward
\emph{earlier} onsets, shrinking the directional means in magnitude by
up to $30$~mN\,m on the noisiest datasets. The route is indirect. At
\emph{fixed} $(C_2, K)$ a baseline biased toward the motion direction
\emph{delays} the detected onset: the post-onset residual
$C_1[\cosh(C_2(t - t_{\mathrm{crit}}^{\mathrm{real}}))
- \cosh(C_2(t - t_{\mathrm{crit}}^{\mathrm{est}}))] - C$ can only vanish
with a later split, and an oppositely biased baseline advances it. The
dominant effect, however, is that the leaked-rise flattening of the cost
landscape moves the CV-selected calibration to a different point on the
$(C_2, K)$ ridge — e.g.\ $C_2\!: 8.00 \to 5.38$, $K\!: 0.060 \to 0.150$
on the most affected dataset — whose net effect is the earlier-onset
magnitude shrink; the ramp-invariance CV itself is essentially
unchanged. The
pivot-free offset $M_{\mathrm{ff}}$ moves by at most $4.2$~mN\,m
($0.13$~mm), well below the run-to-run scatter — the median removes a
systematic bias and tunes nothing. A second ablation pins $C = 0$, the
ideal-sensor value implied by the onset conditions: the measured
pre-onset baseline is in fact nonzero (gyro bias plus residual spin-up
motion; $|C|$ median $3.1$, p90 $14.2$, max $30.6$~mrad/s), yet the
estimate barely moves ($71/132$ onsets identical,
$M_{\mathrm{ff}}$ shift $\le 1.3$~mN\,m). The estimated baseline is
therefore not load-bearing on this dataset; it is kept because it makes
the sweep invariant to the gyro bias — a rig- and temperature-dependent
quantity that need not stay this small — at zero parameter cost. Nor
does $C$ need the bias to be constant in time: it is re-estimated per
run, so only \emph{within-window} drift (windows last $1$--$10$~s)
matters. Measured over the guaranteed pre-onset segment
($|M| < 0.3$~N\,m, spans $\le 3.0$~s), the baseline wanders by
$3.4$~mrad/s (median; p90 $24$, max $68$) — an upper bound that
conflates true sensor drift with real elastic pre-onset motion
(\texttt{analysis/bias\_drift.py}). This wander is present in the very
data on which the $C = 0$ bound and the rate-invariance certificates
were earned, and its structure is covered by the perturbation
absorption: the constant part goes to $C$, a $\tau$-linear trend is
absorbed by the span projection, and only the residual curvature
reaches $\rho$. \emph{Sweep margin.} The margin
$N_{\mathrm{seg}} = 8$, applied at both window ends, guarantees every
candidate split leaves at least $N_{\mathrm{seg}}$ samples in each
segment; it excises only degenerate candidates: near
$j = i_{\mathrm{start}}$ the baseline would be estimated from a handful
of samples, and near $j = i_{\mathrm{end}}$ the post-onset arc is too
short for the cosh branch to be falsifiable (at
$j = i_{\mathrm{end}}$ its cost is nearly an empty sum).
Since physically feasible onsets lie at least
$\hat N_{\mathrm{pre}} \ge 30$ samples inside the window
(Sec.~\ref{sec:gates}) and the window closes only after the response is
well developed, the sweep remains exhaustive over every feasible onset.

\begin{remark}
Because the model curve is fully determined before the post-onset data
are seen, goodness of fit is a \emph{prediction test}, not flexibility:
across all gate-passing runs the normalized post-onset residual is
$\approx 6\%$ (gyro-noise dominated) and invariant over the twelvefold
ramp-rate range — the empirical certificate that the response stays in
the cosh family.
\end{remark}

\section{Rig-constant calibration}
\label{sec:calibration}

Algorithm~\ref{alg:cosh} presupposes $(C_2, K)$. Although both have
geometric expressions \eqref{eq:cosh}, they are \emph{not} computed from
geometry: $z_{\mathrm{CoM}}$ and especially $J_P$ are not precisely known
a priori, and — more importantly — the constants that make
\eqref{eq:cosh} reproduce the measured response are \emph{effective}
values that absorb the state-linear contributions of the unmodeled
effects (the local slope of the gravity nonlinearity and the
ground-effect spring). Both are calibrated once per dataset.

Joint residual minimization is not a usable criterion, and the reason
can be made exact. The identity
$\cosh x - 1 = 2\sinh^2(x/2)$ factors the closed form as
\begin{equation}
C_1\bigl(\cosh(C_2\tau) - 1\bigr) \;=\;
\underbrace{\tfrac{C_1}{2}\,e^{-C_2 t_{\mathrm{on}}}}_{\textstyle A}\;
e^{C_2 t}\,\bigl(1 - e^{-C_2\tau}\bigr)^{2},
\end{equation}
whose last factor rises to one within a few growth times $C_2^{-1}$.
On the tail the amplitude and the onset therefore enter the data only
through the product $A$, and the substitution
\begin{equation}
C_1 \to \lambda C_1, \qquad
t_{\mathrm{on}} \to t_{\mathrm{on}} + C_2^{-1}\ln\lambda
\label{eq:ridge}
\end{equation}
leaves $A$ \emph{exactly} invariant for every $\lambda > 0$: the
response changes only through the last factor, by the relative amount
$\bigl[(1 - \lambda e^{-C_2\tau})/(1 - e^{-C_2\tau})\bigr]^{2} - 1
\approx -2(\lambda - 1)\,e^{-C_2\tau}$, which decays exponentially in
$\tau$. An earlier onset with a smaller amplitude thus trades off
against a later onset with a larger one everywhere except near the
onset — where the response is quadratic, $\tfrac12 C_1 C_2^2\tau^2$,
and of the order of the gyro noise. Equivalently, the sensitivity
vectors $\partial\omega/\partial C_1 \propto \cosh(C_2\tau) - 1$ and
$\partial\omega/\partial t_{\mathrm{on}} \propto \sinh(C_2\tau)$ become
parallel on the tail, so the Gauss--Newton Hessian of the summed
residual is nearly singular along $\delta C_1/C_1 = C_2\,\delta t$ —
a \emph{shallow ridge} in the $(C_2, K)$ plane whose minimizer is set
by noise, not physics. This is an ill-conditioning statement, not a
non-identifiability one: in the noiseless limit the post-onset shape
determines all three parameters uniquely. Under gyro noise, however,
the information along \eqref{eq:ridge} is exponentially small, so
per-run estimates — and with them the detected critical moments —
scatter along the ridge. The calibration should therefore be read as
\emph{variance suppression}: choosing $(C_2, K)$ once per dataset to
minimize the dispersion of $M_{\mathrm{crit}}$ across ramp rates
removes precisely that scatter, which is where the practical strength
of the method lies.

The ridge is invisible to the residual but not to the deliverable.
Moving along \eqref{eq:ridge} shifts every detected onset by the same
\emph{time} offset $\delta t = C_2^{-1}\ln\lambda$ — with
$\lambda = K^\ast/K$ the same for all runs, since $C_1 = K\dot M$
scales them alike — and therefore biases the critical moment by
\begin{equation}
\delta M_{\mathrm{crit}} \;=\; \dot M\,\delta t \;\propto\; \dot M ,
\end{equation}
linearly in the ramp rate: a wrong pair \emph{fans out} the identified
$M_{\mathrm{crit}}$ across the twelvefold rate range, while the correct
pair leaves it rate-invariant. The direction the residual cannot see is
thus exactly the direction a rate-consistency criterion sees best. We
therefore calibrate by \emph{physical self-consistency}: a static
threshold must not depend on how fast it was approached. With runs at
ramp rates
$r \in \{0.10,\dots,1.20\}$~N\,m/s,
\begin{equation}
\mathrm{score}(C_2, K) \;=\; \sum_{\mathrm{dir}\in\{+,-\}}
\mathrm{CV}\Bigl(\bigl\{ M_{\mathrm{crit}}^{(r)}(C_2, K)\bigr\}_{r}\Bigr),
\qquad \mathrm{CV} = \mathrm{std}/|\mathrm{mean}| ,
\label{eq:score}
\end{equation}
where $M_{\mathrm{crit}}^{(r)}$ is the output of
Algorithm~\ref{alg:cosh} on the run with commanded rate $r$.

\paragraph{Search ranges.} The admissible box is set so that it covers
every physically plausible constant with margin. For
$K = 1/(W z_{\mathrm{CoM}})$, the coarse range
$K \in [0.10, 0.50]$~(N\,m)$^{-1}$ corresponds to
$z_{\mathrm{CoM}} \in [0.06, 0.32]$~m at the vehicle weight — bracketing
the plausible CoM heights — with the refinement allowed to move within
the wider clip box $[0.05, 0.70]$ ($z_{\mathrm{CoM}}$ up to
$0.63$~m). For $C_2 \in [3, 8]$~rad/s, the implied pivot inertia
$J_P = 1/(K C_2^2)$ spans roughly $[0.03, 1.1]$~kg\,m$^2$, bracketing
the parallel-axis estimate
$J + m z_{\mathrm{CoM}}^2 \approx 0.2$--$0.35$~kg\,m$^2$ by a factor of
several on both sides.

\begin{algorithm}[h]
\caption{Ramp-invariance consistency score}
\label{alg:score}
\begin{algorithmic}[1]
\Require gated runs at rates $r \in \{0.10,\dots,1.20\}$~N\,m/s, both
  tip directions; candidate pair $(C_2, K)$
\Ensure consistency score $s$ \Comment{lower = $M_{\mathrm{crit}}$ more
  rate-invariant}
\For{each gated run $r$}
  \State $M_{\mathrm{crit}}^{(r)} \gets$
    Algorithm~\ref{alg:cosh} with $(C_2, K)$ pinned
    \Comment{stride-3 sweep}
\EndFor
\For{$\mathrm{dir} \in \{+,-\}$}
  \State $\mathrm{CV}_{\mathrm{dir}} \gets
    \mathrm{std}\bigl(\{M_{\mathrm{crit}}^{(r)}\}_{r\in\mathrm{dir}}\bigr)
    \,/\, \bigl|\mathrm{mean}\bigl(\{M_{\mathrm{crit}}^{(r)}\}_{r\in
    \mathrm{dir}}\bigr)\bigr|$
\EndFor
\State \Return $s \gets \mathrm{CV}_{+} + \mathrm{CV}_{-}$
  \Comment{Eq.~\eqref{eq:score}: a static threshold must not depend on
    the approach rate}
\end{algorithmic}
\end{algorithm}

\begin{algorithm}[h]
\caption{Rig-constant calibration (coarse-to-fine)}
\label{alg:calib}
\begin{algorithmic}[1]
\Require gated runs; box $C_2 \in [3,8]$~rad/s,
  $K \in [0.05,0.70]$~(N\,m)$^{-1}$
\Ensure calibrated pair $(C_2^\ast, K^\ast)$
\State \textbf{coarse:} evaluate Algorithm~\ref{alg:score} on the
  $11{\times}11$ grid $C_2 = 3.0\!:\!0.5\!:\!8.0$,
  $K = 0.10\!:\!0.04\!:\!0.50$
\State $(C_2^{b}, K^{b}) \gets \arg\min$ over the coarse grid
\State \textbf{fine:} evaluate Algorithm~\ref{alg:score} on a
  $9{\times}9$ uniform grid over
  $[C_2^{b} \pm 0.5] \times [K^{b} \pm 0.04]$, clipped to the box
  \Comment{resolution $0.125$~rad/s, $0.01$~(N\,m)$^{-1}$}
\State $(C_2^\ast, K^\ast) \gets \arg\min$ over the fine grid
\end{algorithmic}
\end{algorithm}

\paragraph{Coarse-to-fine procedure.}
Algorithm~\ref{alg:score} states the consistency criterion as a
stand-alone routine — it wraps Algorithm~\ref{alg:cosh} and returns the
directional sum of coefficients of variation — and
Algorithm~\ref{alg:calib} performs the search over it. The coarse scan
evaluates
\eqref{eq:score} on an $11 \times 11$ grid
($\Delta C_2 = 0.5$~rad/s, $\Delta K = 0.04$~(N\,m)$^{-1}$;
$121$ evaluations). A single refinement pass then re-scans a
$9 \times 9$ uniform grid spanning $\pm$ one coarse step about the
coarse minimizer ($\pm 0.5$~rad/s, $\pm 0.04$~(N\,m)$^{-1}$, clipped to
the box), i.e.\ final resolutions of $0.125$~rad/s and
$0.01$~(N\,m)$^{-1}$ — so the reported constants are not quantized to
the coarse step. One pass suffices because the objective is flat along
the $(C_2, K)$ ridge; it is also \emph{piecewise constant} (the onset
is an integer index, so small parameter changes leave
$M_{\mathrm{crit}}$ unchanged), which is why gradient-based or simplex
optimizers stall on it, and a population-based global search
(differential evolution) was found to cost about $4\times$ more for no
measurable gain. During calibration the onset sweep of
Algorithm~\ref{alg:cosh} is subsampled by a stride of $3$ samples for
speed; the final identification runs at full resolution with the
selected pair. Each evaluation is an arithmetic sweep, so no iterative
optimizer is involved anywhere in the pipeline.

\begin{remark}
The objective \eqref{eq:score} retains a shallow valley, so the
individual constants are only loosely determined:
$W z_{\mathrm{CoM}} = 1/K$ and $J_P = 1/(K C_2^2)$ are
order-of-magnitude sanity checks, not precision measurements. What the
criterion pins down — by construction — is the identified onset:
$M_{\mathrm{crit}}$ is robust along the valley, and a $\pm 20\%$
mis-specification of either constant moves the identified CoM by less
than $0.9$~mm.
\end{remark}

\section{Aggregation and deliverables}
\label{sec:aggregation}

Directional means over the gated runs give $\bar M_{+}$ and
$\bar M_{-}$. The take-off feed-forward moment is the \emph{pivot-free
average}
\begin{equation}
M_{\mathrm{ff}} = \tfrac12\,(\bar M_{+} + \bar M_{-}) ,
\label{eq:mff}
\end{equation}
in which the pivot-arm terms $\pm(W - f)\,l_p$ — and any
thrust-proportional contamination acting through the arm $l_p$, ground
effect included — cancel by antisymmetry, while landing-gear asymmetry
is deliberately \emph{retained}: the same contact-phase moment acts at
lift-off, where $M_{\mathrm{ff}}$ compensates it. The CoM offset follows
from the moment-to-offset conversion of $M_{\mathrm{ff}}$ and is used
for validation against load-cell ground truth only; $M_{\mathrm{ff}}$
itself is the quantity applied in flight. A roll excitation ($M_x$)
senses the $y$-offset via $M_{\mathrm{ff},x} = +W y_{\mathrm{off}}$; a
pitch excitation ($M_y$) senses the $x$-offset via
$M_{\mathrm{ff},y} = -W x_{\mathrm{off}}$.

\section{Comparison with per-run nonlinear least squares}
\label{sec:nls}

The predecessor of Algorithm~\ref{alg:cosh} fitted each run
individually by nonlinear least squares (\textsc{scipy}
\texttt{least\_squares}, trust-region reflective): free
$(C_1, C_2, C)$ with $C_2$ bounded to the physical band
$[3, 8]$~rad/s and a free baseline, the onset swept locally around a
quadratic-model seed. To quantify what the dataset-coupled calibration
buys, both estimators were run on \emph{identical} inputs — same
signals, excitation windows, allocator caps and onset-free gates (the
gated run set is method-independent because the gates use the moment
trace only), $132$ runs in total
(\texttt{analysis/nls\_comparison.py}). Table~\ref{tab:nls} compares
the deliverables — the pivot-free offset \eqref{eq:mff} and the CoM
offset — against the load-cell ground truth;
Table~\ref{tab:nls-disp} compares the directional critical-moment
means and their ramp-rate dispersion.

\begin{table}[t]
\centering
\small
\setlength{\tabcolsep}{5pt}
\begin{tabular}{cc r rrr rrr}
\toprule
 & & $\lambda_{\mathrm{truth}}$ &
\multicolumn{3}{c}{COSH} & \multicolumn{3}{c}{NLS} \\
\cmidrule(lr){4-6} \cmidrule(lr){7-9}
case & axis & [mm] & $M_{\mathrm{ff}}$ [N\,m] & $\hat\lambda$ [mm] &
$|\mathrm{err}|$ & $M_{\mathrm{ff}}$ [N\,m] & $\hat\lambda$ [mm] &
$|\mathrm{err}|$ \\
\midrule
1 & x & $-2.90$  & $-0.052$ & $-1.73$  & $1.17$ & $-0.062$ & $-2.05$  & $0.85$ \\
1 & y & $-11.45$ & $+0.329$ & $-10.93$ & $0.52$ & $+0.328$ & $-10.91$ & $0.54$ \\
2 & x & $-14.29$ & $-0.481$ & $-15.22$ & $0.93$ & $-0.472$ & $-14.95$ & $0.66$ \\
2 & y & $-9.90$  & $+0.327$ & $-10.36$ & $0.46$ & $+0.306$ & $-9.67$  & $0.23$ \\
3 & x & $-5.26$  & $-0.272$ & $-8.62$  & $3.36$ & $-0.270$ & $-8.54$  & $3.28$ \\
3 & y & $+3.14$  & $-0.023$ & $+0.73$  & $2.41$ & $-0.033$ & $+1.06$  & $2.08$ \\
4 & x & $+6.67$  & $+0.166$ & $+5.26$  & $1.41$ & $+0.156$ & $+4.94$  & $1.73$ \\
4 & y & $+2.40$  & $-0.005$ & $+0.16$  & $2.24$ & $+0.020$ & $-0.64$  & $3.04$ \\
5 & x & $+10.91$ & $+0.338$ & $+10.70$ & $0.21$ & $+0.317$ & $+10.05$ & $0.86$ \\
5 & y & $-10.89$ & $+0.338$ & $-10.70$ & $0.19$ & $+0.332$ & $-10.50$ & $0.39$ \\
\midrule
\multicolumn{3}{r}{median / RMS / max $|\mathrm{err}|$ [mm]} &
\multicolumn{3}{c}{$1.05$ / $1.64$ / $3.36$} &
\multicolumn{3}{c}{$0.86$ / $1.72$ / $3.28$} \\
\bottomrule
\end{tabular}
\caption{Deliverables against the load-cell ground truth
(manuscript Table~7): pivot-free offset \eqref{eq:mff} and the CoM
offset $\hat\lambda$, per-case weight ($3.066$~kg for case~1,
$3.220$~kg otherwise). One row per case--axis; the two estimators sit
side by side.}
\label{tab:nls}
\end{table}

\begin{table}[t]
\centering
\small
\setlength{\tabcolsep}{5pt}
\begin{tabular}{cc rrrr rrrr}
\toprule
 & & \multicolumn{4}{c}{COSH} & \multicolumn{4}{c}{NLS} \\
\cmidrule(lr){3-6} \cmidrule(lr){7-10}
case & axis & $\bar M_{-}$ & $\bar M_{+}$ & $\mathrm{CV}_{-}$ &
$\mathrm{CV}_{+}$ & $\bar M_{-}$ & $\bar M_{+}$ & $\mathrm{CV}_{-}$ &
$\mathrm{CV}_{+}$ \\
 & & [N\,m] & [N\,m] & [\%] & [\%] & [N\,m] & [N\,m] & [\%] & [\%] \\
\midrule
1 & x & $-1.054$ & $0.950$ & $\mathbf{3.8}$ & $\mathbf{5.7}$ & $-1.018$ & $0.894$ & $6.4$ & $7.7$ \\
1 & y & $-0.422$ & $1.079$ & $\mathbf{2.4}$ & $\mathbf{2.2}$ & $-0.399$ & $1.056$ & $8.8$ & $7.2$ \\
2 & x & $-1.829$ & $0.867$ & $\mathbf{1.2}$ & $\mathbf{1.5}$ & $-1.792$ & $0.848$ & $2.5$ & $4.2$ \\
2 & y & $-0.666$ & $1.320$ & $\mathbf{1.2}$ & $\mathbf{1.6}$ & $-0.637$ & $1.248$ & $4.9$ & $5.8$ \\
3 & x & $-1.542$ & $0.997$ & $\mathbf{1.7}$ & $\mathbf{2.4}$ & $-1.573$ & $1.033$ & $2.4$ & $4.1$ \\
3 & y & $-0.983$ & $0.937$ & $\mathbf{1.7}$ & $\mathbf{2.6}$ & $-0.990$ & $0.923$ & $3.2$ & $7.0$ \\
4 & x & $-1.136$ & $1.468$ & $\mathbf{3.4}$ & $\mathbf{2.2}$ & $-1.151$ & $1.463$ & $3.8$ & $3.5$ \\
4 & y & $-0.998$ & $0.988$ & $\mathbf{2.9}$ & $\mathbf{3.7}$ & $-0.928$ & $0.969$ & $10.3$ & $6.1$ \\
5 & x & $-1.019$ & $1.695$ & $\mathbf{3.9}$ & $\mathbf{2.1}$ & $-0.989$ & $1.624$ & $7.2$ & $2.6$ \\
5 & y & $-0.635$ & $1.310$ & $\mathbf{6.5}$ & $\mathbf{2.9}$ & $-0.577$ & $1.241$ & $15.2$ & $5.4$ \\
\bottomrule
\end{tabular}
\caption{Directional critical-moment means and their ramp-rate
dispersion. Bold marks the lower CV of each pair — COSH in all
$20/20$ groups.}
\label{tab:nls-disp}
\end{table}

Three observations. \emph{(i) Run-level robustness is where the
methods separate:} the directional CV is lower under COSH in
$20/20$ case--axis--direction groups — median $2.4\%$ vs.\ $5.6\%$,
NLS worst case $15.2\%$ — and the paired per-run critical moments
differ by $27.6$~mN\,m at the median but $118$~mN\,m at the 90th
percentile and up to $217$~mN\,m. This is the ridge of
Sec.~\ref{sec:calibration} realized: the free fit scatters along the
amplitude--onset trade-off run by run, and its directional magnitudes
sit systematically low ($\lvert\bar M_\pm\rvert$ smaller in $17/20$
groups) — the early-onset drift of the unconstrained fit. \emph{(ii)
The deliverable survives either way:} against the load-cell truth the
CoM error is statistically indistinguishable (COSH median
$1.05$ / RMS $1.64$ / max $3.36$~mm; NLS $0.86$ / $1.72$ / $3.28$~mm),
because the pivot-free average \eqref{eq:mff} cancels the largely
common-mode early-onset bias between the $+$ and $-$ directions, and
averaging $5$--$7$ runs per direction suppresses the scatter.
\emph{(iii) The calibration is what converts per-run estimates from
indicative to trustworthy:} with COSH a \emph{single} run already sits
within a few percent of the directional mean, which is what enables
run-count reduction and per-run gating diagnostics; the NLS estimator
needs the full ensemble to average its ridge scatter away.

\section{Reproducibility}

Every number in this document is reproducible from the repository:
\texttt{critical\_value\_getter\_piecewise.py} (pipeline, gates, caps,
calibration, Algorithm~\ref{alg:cosh});
\texttt{analysis/ramp\_quality.py} (gates);
\texttt{analysis/sensitivity.py} ($\pm20\%$ robustness);
\texttt{analysis/cosh\_fidelity.py} (prediction-test residuals);
\texttt{analysis/excitation\_angle\_design.py} (rate/cutoff design);
\texttt{analysis/ge\_linearity.py}, \texttt{ge\_trajectory.py}
(ground-effect bracket); \texttt{analysis/tiltcap\_ablation.py} (window
ablation); \texttt{analysis/baseline\_stat\_ablation.py} (median/mean/%
zero baseline); \texttt{analysis/bias\_drift.py} (in-window baseline
drift); \texttt{analysis/nls\_comparison.py}
(Table~\ref{tab:nls}).

\end{document}
